% Generated draft: requires mathematical and source review.
\documentclass[11pt]{article}
\usepackage[margin=25mm]{geometry}
\usepackage{fontspec}
\setmainfont{Latin Modern Roman}
\usepackage{amsmath,amssymb,mathtools,unicode-math}
\setmathfont{Latin Modern Math}
\usepackage{enumitem,needspace,xurl,hyperref}
\setlength{\parindent}{0pt}
\setlength{\parskip}{5pt}
\setlength{\emergencystretch}{3em}
\newcommand{\N}{\mathbb{N}}
\newcommand{\Z}{\mathbb{Z}}
\newcommand{\Q}{\mathbb{Q}}
\newcommand{\R}{\mathbb{R}}
\newcommand{\C}{\mathbb{C}}
\newcommand{\F}{\mathbb{F}}
\newcommand{\A}{\mathbb{A}}
\newcommand{\PP}{\mathbb{P}}
\newcommand{\E}{\mathbb{E}}
\newcommand{\eps}{\varepsilon}
\newcommand{\dd}{\,\mathrm{d}}
\begin{document}
\section*{UnsolvedMath ID: 30005467}
\textbf{Source problem code:} catalogue entry
\subsection{Definitions and mathematical statement}
For a positive integer \(n\), let \([n]=\{1,\ldots,n\}\). For \(A\in\C^{n\times n}\) and \(S\subseteq[n]\), let \(A[S,S]\) denote the principal submatrix indexed by \(S\), and put \(p_S(A)=\det A[S,S]\), with \(p_{\varnothing}(A)=1\). Define the principal-minor map by
\[
\Phi_n:\C^{n\times n}\longrightarrow H_n:=\{(x_S)_{S\subseteq[n]}\in\C^{2^n}:x_{\varnothing}=1\},\qquad A\longmapsto(p_S(A))_{S\subseteq[n]}.
\]
Including the empty principal minor is a coordinate convention; deleting that fixed coordinate identifies \(H_n\) with \(\C^{2^n-1}\). The matrices here are arbitrary complex square matrices, with no symmetry or Hermitian assumption.

The first question is to determine equations describing \(\Phi_5(\C^{5\times5})\). Under the literal interpretation as an exact description by polynomial equations, this asks for a family
\[
\mathcal F_5\subseteq\C[x_S:\varnothing\ne S\subseteq[5]]
\]
such that, for every \(x\in H_5\),
\[
\bigl(\exists A\in\C^{5\times5}\;\forall S\subseteq[5],\ x_S=\det A[S,S]\bigr)
\quad\Longleftrightarrow\quad
\bigl(\forall f\in\mathcal F_5,\ f((x_S)_{S\ne\varnothing})=0\bigr).
\]
The existence of such an exact description is not assumed. The supplied record does not specify whether ``image'' instead means its Zariski closure; equations defining that closure would be a different target unless equality with the image is established.

The second question asks whether, for general positive integers \(n\), equations describing the principal-minor image can be obtained using
\[
G_n=\operatorname{SL}_2(\C)^n\rtimes S_n,
\]
where \(S_n\) permutes the \(n\) factors, together with what the original passage calls ``some increasing combinatorial structures.'' The cleaned statement calls these ``compatible combinatorial structures.'' The record supplies neither the relevant representation or ambient space for the group action, nor the initial equations, nor a definition of these structures or their compatibility. Thus this part remains a qualified research question, rather than a fully specified equation-generation assertion. In particular, an affine, homogeneous, or projective interpretation has not been selected.

The original passage also asks ``Does the converse hold?'' Its antecedent is absent from the supplied record, so that converse cannot be formulated precisely here.
\subsection{Short English statement}
Determine equations describing the collections of principal minors of arbitrary complex \(5\times5\) matrices, and investigate a construction for general matrix size using \(\operatorname{SL}_2(\C)^n\rtimes S_n\) and additional combinatorial structures. The record leaves unresolved whether the intended target is the exact image or its Zariski closure, how the group acts, which combinatorial structures are intended, and which converse is being asked about.
\subsection{Sources}
\begin{itemize}
\item Local UnsolvedMath registry, permanent ID \(30005467\), title: \emph{Images of Principal Minor Maps for Square Matrices}; source problem code \texttt{catalogue entry}. Registry locations identified by the project are \texttt{data/unsolvedmath/problems.json} and \texttt{catalog/problems.json}; this draft uses the supplied complete record. The proposed year is 2023; no individual proposer is named.
\item \emph{Real Algebraic Geometry with a View toward Koopman Operator Methods} (2023), Oberwolfach Reports, pp.\ 741--811, DOI \texttt{10.4171/owr/2023/14}. Supplied report reference: \url{https://doi.org/10.4171/owr/2023/14}. The record describes Oberwolfach Reports as documenting mathematical workshops and their open-problem discussions; no specific page within the report is supplied for this problem.
\item Al Ahmadieh--Vinzant, ``symmetric principal minors'' (supplied reference label): \url{https://arxiv.org/abs/2105.13444}.
\item Al Ahmadieh--Vinzant, ``determinantal representations'' (supplied reference label): \url{https://arxiv.org/abs/2205.05267}.
\item Supplied literature assessment, checked \texttt{2026-08-22}: the registry labels the problem \texttt{partially\_solved}. Its assessment states that the complex symmetric principal-minor problem and its invariant equations are characterized, and that Hermitian images have finite equation-and-inequality descriptions, while the general square-matrix image and the full bundled converse questions remain unresolved. This is supplied dated curation derived from the OWR curation, not an independently verified literature assessment; its status and cited claims require independent verification.
\item The registry labels the statement \texttt{corrected\_verified} and reports that the official report was checked directly and that the cleaned formulation restores definitions and removes extraction artifacts or source-dependent wording. Those are supplied verification claims; this draft does not independently verify them or claim to have visited any source.
\end{itemize}
\end{document}
